Phishing defenders need a verdict they can trust and a reason they can check. Large
language models are accurate but score trust at the label, so several models can agree on
the right answer while citing different or fabricated evidence. We propose \toolname{}, a
selective detector that reads trust from grounded cross-model evidence agreement (GEA): a
diverse panel emits typed cues from the page's URL, DOM, screenshot, and logo (certificate
and redirect cues when available), verification tools confirm them against the page, and a
held-out split calibrates GEA to a probability of correctness. \toolname{} acts only on
grounded agreement and otherwise abstains, returning verified cues as an explanation.
On \textsc{PhishSel} ($998$ test pages of feed-sourced phishing paired with legitimate login
pages), \toolname{} attains the best-calibrated trust (ECE $0.015$, statistically
significant) and the highest selective accuracy at $80\%$ coverage ($97.0\%$, significant).
Its area-under-risk--coverage is lowest in point estimate ($2.3$) though not statistically
separated from the strongest baseline; it trades top-end deployable coverage (Cov$_{99}$
$0.25$ vs.\ $0.36$) for sharper calibration in the $50$--$80\%$ deployment band. Under
evidence-targeted attacks, abstention on evaded pages reaches $97.5$--$100\%$ while
single-model confidence stays maximally certain; even a white-box, all-cue
adversary -- cloaking the form action, erasing the brand's textual identity, and
occluding the logo simultaneously -- flips more verdicts but \toolname{} abstains
on every one of them ($100\%$, $n{=}50$). Calibrated trust transfers to
unseen brands, and negative ablations show that learned aggregation, larger text models, and
cheaper cascades do not beat the parameter-free design.
